\documentclass[11pt]{article}

\usepackage[margin=1in]{geometry}
\usepackage{parskip}
\usepackage{amsmath, amssymb, amsthm}
\usepackage{booktabs}
\usepackage{hyperref}

\title{A small library for simulating and analysing sums of fair dice}
\author{Group N}
\date{}

\begin{document}
\maketitle

\begin{abstract}
We present \texttt{dice}, a small Python library for simulating sums
of fair \(s\)-sided dice and comparing the empirical distribution
with the exact theoretical distribution. We use the standard library
throughout: \texttt{random} for the simulation, \texttt{statistics}
for the summary numbers, and \texttt{fractions.Fraction} for the
exact distribution. We recover the textbook values
\(P(s = 2) = 1/36\), \(P(s = 7) = 1/6\) and \(P(s = 12) = 1/36\) for
two six-sided dice, and we show that the empirical distribution from
\(10000\) simulated rolls agrees with the theoretical values to
within sampling error. The library follows the Diataxis framework for
documentation~\cite{procida2017} and the conventions of the
open-source textbook \emph{Python for
Mathematics}~\cite{knight2024}.
\end{abstract}

\section{Introduction}

The distribution of the sum of \(n\) fair \(s\)-sided dice is a
standard probability example. For small \(n\) the distribution can be
written out by hand: for \(n = 2\) and \(s = 6\) the well-known
triangular distribution has \(P(s = 2) = 1/36\), \(P(s = 7) = 6/36\),
and the rest of the values fall symmetrically between. For larger
\(n\) the closed-form expression in terms of binomial coefficients is
awkward to state and easy to get wrong.

We wanted a small library that lets us both simulate the sum and
compute the exact distribution, with the latter as a sanity check on
the former. The standard library provides everything we need:
\texttt{random} for the simulation, \texttt{statistics} for the
summary numbers, and \texttt{fractions.Fraction} for the rationals.

We present a five-function library,
\texttt{dice}; we recover the textbook values for two six-sided dice;
and we show that the empirical distribution from \(10000\) rolls is
within a comfortable margin of the theoretical values.

\section{The library}

The library lives in the single module \texttt{dice.py} and follows
the probability chapter of \emph{Python for
Mathematics}~\cite{knight2024}.

\texttt{roll(number\_of\_dice, sides, seed)} returns
a tuple of face values, each drawn uniformly from
\(\{1, 2, \dots, s\}\). An optional ``seed`` argument lets us
reproduce a roll.

\texttt{simulate\_sums(number\_of\_dice,
number\_of\_rolls, sides, seed)} returns a list of sums, one per roll.

\texttt{empirical\_distribution(samples)} returns a dictionary
mapping each observed value to its relative frequency.

\texttt{theoretical\_distribution(number\_of\_dice, sides)} computes
the exact distribution by repeated convolution of the single-die
distribution. We use \texttt{fractions.Fraction} so that the
probabilities are exact rationals.

\texttt{summary\_statistics(samples)}
returns the mean, median, and sample standard deviation, using
\texttt{statistics.mean}, \texttt{statistics.median} and
\texttt{statistics.stdev}.

\section{Worked example}

We illustrate the library on the sum of two six-sided dice. The
theoretical distribution is the triangular distribution shown in
Table~\ref{tab:two-d6}. We then simulate \(10000\) rolls with
seed \(0\); the empirical probability of rolling a \(7\) is
\(0.16{\sim}0.17\), which lies inside the interval \([0.13, 0.20]\)
that we use in the doctest.

\begin{table}[h]
\centering
\begin{tabular}{l l l l l l l l l l l l}
\toprule
Sum & 2 & 3 & 4 & 5 & 6 & 7 & 8 & 9 & 10 & 11 & 12 \\
Prob.\ ($\times 36$) & 1 & 2 & 3 & 4 & 5 & 6 & 5 & 4 & 3 & 2 & 1 \\
\bottomrule
\end{tabular}
\caption{\textbf{The exact distribution of the sum of two six-sided
dice.} The probability of each total is shown as a multiple of
\(1/36\). The mode is at \(7\), with probability \(6/36 = 1/6\).
The library returns these values exactly using
\texttt{fractions.Fraction}.}
\label{tab:two-d6}
\end{table}

\section{Discussion}

The probability chapter of \emph{Python for Mathematics} reaches for
\texttt{sympy.Rational}, which would also work. We chose
\texttt{fractions.Fraction} because the rest of the library is
standard-library only, and bringing in \texttt{sympy} would be
disproportionate for the small amount of exact arithmetic we need.

We only handle fair dice. A natural extension
would be to accept a custom probability mass function on the faces of
each die, so that we can model loaded dice or dice with non-standard
faces. The convolution code already accepts a dictionary, so this
extension would be small.

\section{Conclusion}

We have presented \texttt{dice}, a small library for simulating sums
of fair dice and comparing them against the exact distribution. The
library recovers the textbook values for two six-sided dice and the
empirical distribution from a seeded simulation agrees with them to
within sampling error.

\bibliographystyle{plain}
\bibliography{references}

\end{document}
